\chapter{System Architecture \& Infrastructure at Scale}
\label{system-architecture-infrastructure-at-scale}

Training LLMs with reinforcement learning from human feedback is as much a systems engineering challenge as it is an algorithmic one. Unlike standard supervised fine-tuning---which involves a single model, a single forward-backward pass, and well-understood scaling---RLHF requires \emph{multiple models} (policy, reference, reward model, value head) to be loaded simultaneously, coordinated through a complex rollout-scoring-training loop, and distributed across dozens to hundreds of GPUs. This chapter covers the systems-level details that make large-scale RLHF training possible: memory budgeting, parallelism strategies (Data, Tensor, Pipeline, Sequence, and their combinations), the generation bottleneck, decoupled architectures, weight synchronization, fault tolerance, and production monitoring.

\section{The 4-Model Memory Challenge}
\label{the-4-model-memory-challenge}

\begin{figure}[ht!]
\centering
\includegraphics[width=0.85\textwidth]{figures/fig_031_fig31.png}
\caption{70B PPO memory budget: the four models required for RLHF and their memory footprints. Total: 1470--1560GB. Minimum 19--20 A100-80GB (naive). With ZeRO-3: fits in 8 nodes.}
\end{figure}

\begin{warningbox}[Memory Budget Reality Check -- 70B BF16]
\begin{tabular}{@{}ll@{}}
\toprule
Policy weights (BF16) & 140 GB \\
\midrule
FP32 master weights & 280 GB \\
Adam optimizer (m + v, FP32) & 560 GB \\
Gradients (BF16) & 140 GB \\
Reference model & 140 GB (or 70 GB in INT8) \\
Reward model & 140 GB (or 70 GB in INT8) \\
Activations (batch 128, seq 2048) & 50--100 GB \\
KV cache for generation & 20--60 GB \\
\textbf{Total} & \textbf{1470--1560 GB} \\
\bottomrule
\end{tabular}

\smallskip
$\div$ 80 GB/GPU = \textbf{19--20 A100s minimum} (without any parallelism overhead).
\end{warningbox}

\section{Parallelism Strategies in Detail}
\label{sec:parallelism}

Training large language models requires distributing computation across many GPUs. There are fundamentally different \emph{axes} along which to parallelize, each with distinct trade-offs. This section provides detailed coverage of each strategy with mathematical formulations, diagrams, and practical guidance.

\begin{figure}[ht!]
\centering
\includegraphics[width=0.85\textwidth]{figures/fig_032_fig32.png}
\caption{Overview of the four parallelism strategies. Production systems typically combine 2--3 of these simultaneously.}
\end{figure}

\subsection{Data Parallelism (DP) and Distributed Data Parallelism (DDP)}
\label{subsec:dp-ddp}

Data Parallelism is the simplest and most common form of distributed training~\cite{li2020pytorch}. Each GPU holds a \emph{complete copy} of the model, processes a different mini-batch, and synchronizes gradients.

\paragraph{Vanilla DP (PyTorch \texttt{DataParallel}).}
\label{vanilla-dp-pytorch-dataparallel.}

A single-process approach where one ``master'' GPU scatters input, gathers outputs, and broadcasts gradients. Limited by GIL and PCIe bandwidth to the master GPU.

\paragraph{Distributed Data Parallelism (DDP, \texttt{DistributedDataParallel}).}
\label{distributed-data-parallelism-ddp-distributeddataparallel.}

Multi-process: each GPU runs its own process. Gradients are synchronized via ring-AllReduce~\cite{sergeev2018horovod} in the background while backward computation continues.

\begin{figure}[ht!]
\centering
\includegraphics[width=0.85\textwidth]{figures/fig_033_ddp.png}
\caption{DDP: each GPU holds a full model replica and processes a different batch. Gradients are averaged via ring AllReduce, overlapped with backward computation.}
\label{fig:ddp}
\end{figure}

\textbf{Key properties of DDP}:

\begin{itemize}
  \item \textbf{Memory}: Each GPU stores full model + optimizer + gradients. For 70B BF16: $\sim$560~GB/GPU---impossible without memory optimizations.
  \item \textbf{Communication}: One AllReduce of gradient tensor per step. Size = model parameters $\times$ 2 bytes (BF16). Ring AllReduce cost: $2 \cdot \frac{N-1}{N} \cdot M$ bytes transferred per GPU.
  \item \textbf{Scaling}: Near-linear up to $\sim$64 GPUs. Beyond that, communication starts to dominate.
  \item \textbf{Gradient bucketing}: DDP groups parameters into buckets (default 25~MB) and starts AllReduce as soon as a bucket’s gradients are ready---overlapping communication with backward computation.
\end{itemize}

\begin{lstlisting}[style=pythonstyle]
import torch.distributed as dist
from torch.nn.parallel import DistributedDataParallel as DDP

dist.init_process_group(backend="nccl")  # NCCL for GPU communication
model = model.to(local_rank)
model = DDP(model, device_ids=[local_rank],
            gradient_as_bucket_view=True,    # Memory optimization
            static_graph=True)               # Enable comm optimizations
\end{lstlisting}

\begin{warningbox}[DP vs DDP — Always Use DDP]
PyTorch’s legacy \texttt{DataParallel} (DP) should \textbf{never} be used for LLM training:

\begin{itemize}
  \item Single-process, limited by Python GIL
  \item All gradients funnel through GPU 0 (bottleneck)
  \item 2--3$\times$ slower than DDP even on a single node
  \item Cannot scale beyond one machine
\end{itemize}

DDP is the \emph{minimum} parallelism strategy. For LLMs $>$7B, FSDP/ZeRO is preferred.
\end{warningbox}

\subsection{Tensor Parallelism (TP)}
\label{subsec:tensor-parallel}

Tensor Parallelism (Megatron-LM style~\cite{shoeybi2019megatron}) splits \emph{individual weight matrices} across GPUs. Each GPU computes a partial result, and an AllReduce combines them.

\paragraph{Column-Parallel Linear Layer.}
\label{column-parallel-linear-layer.}

The weight matrix $W \in \mathbb{R}^{d \times h}$ is split column-wise across $T$ GPUs: 
\begin{equation}
W = [W_0 \;|\; W_1 \;|\; \cdots \;|\; W_{T-1}], \quad W_i \in \mathbb{R}^{d \times h/T}
\end{equation}
 Each GPU $i$ computes $Y_i = XW_i$ independently (no communication). The output is split along the hidden dimension.

\paragraph{Row-Parallel Linear Layer.}
\label{row-parallel-linear-layer.}

The weight matrix is split row-wise: $W = [W_0; W_1; \ldots; W_{T-1}]$ where $W_i \in \mathbb{R}^{d/T \times h}$. Input $X$ must also be split. Each GPU computes a partial sum, then an \textbf{AllReduce} produces the final output.

\begin{figure}[ht!]
\centering
\includegraphics[width=0.85\textwidth]{figures/fig_034_tp-column.png}
\caption{Column-parallel linear layer (TP=2). The weight is split column-wise; each GPU computes $XW_i$ independently. The MLP pairs this with a row-parallel layer to avoid redundant AllReduce.}
\end{figure}

\paragraph{Transformer Block with TP.}
\label{transformer-block-with-tp.}

In a Transformer layer, Megatron-LM applies TP as follows:

\begin{enumerate}
  \item \textbf{MLP}: Column-parallel for the first linear ($h \to 4h$), row-parallel for the second ($4h \to h$). One AllReduce after the row-parallel layer.
  \item \textbf{Attention}: $Q$, $K$, $V$ projections are column-parallel (split heads across GPUs). Output projection is row-parallel. One AllReduce after output projection.
  \item \textbf{Total}: 2 AllReduce per transformer layer (one for attention, one for MLP).
\end{enumerate}

\begin{figure}[ht!]
\centering
\includegraphics[width=0.85\textwidth]{figures/fig_035_tp-transformer.png}
\caption{Tensor Parallel communication pattern in one Transformer block. Two AllReduce operations (marked in red) are required per layer---one after attention, one after MLP.}
\label{fig:tp-transformer}
\end{figure}

\begin{intuitionbox}[Why TP is Restricted to Intra-Node]
Each transformer layer requires 2 AllReduce operations (marked as $f$ and $g$ above). For a 70B model with 80 layers, that’s 160 AllReduce operations per forward pass (320 including backward). At NVLink speeds (600~GB/s), each AllReduce takes $<$0.5~ms. But over InfiniBand (50~GB/s), the same operation takes $\sim$4~ms, making the total overhead 160 $\times$ 4 = 640~ms---\emph{longer than the computation itself}.

\textbf{Rule}: TP degree $\leq$ GPUs per node (typically TP $\leq$ 8). Use DP/FSDP for inter-node scaling.
\end{intuitionbox}

\begin{keybox}[TP Degree Selection]
\begin{itemize}
  \item \textbf{TP=1}: No tensor parallelism. Model fits on one GPU (typically $\leq$ 13B with BF16).
  \item \textbf{TP=2}: Minimal split. Good for 13--34B inference on 2 GPUs. Low overhead ($<$5\%).
  \item \textbf{TP=4}: Standard for 34--70B inference. Overhead 8--12\%.
  \item \textbf{TP=8}: Full node. Required for 70B+ training. Overhead 12--18\%.
  \item \textbf{TP$>$8}: Cross-node TP. Rarely used---only for 200B+ models where PP alone is insufficient. Overhead 30--50\%.
\end{itemize}

\textbf{Important}: Number of attention heads must be divisible by TP degree. For LLaMA-70B (64 heads), valid TP = 1, 2, 4, 8, 16, 32, 64.
\end{keybox}

\subsection{Sequence Parallelism (SP)}
\label{subsec:sequence-parallel}

Sequence Parallelism~\cite{korthikanti2023reducing} addresses a memory bottleneck that Tensor Parallelism alone cannot solve: the \textbf{activation memory} in LayerNorm and Dropout layers.

\paragraph{The Problem.}
\label{the-problem.}

With TP, weight memory is split across GPUs. But LayerNorm and Dropout operate on the \emph{full} hidden dimension and are replicated on every GPU. Their activations (needed for backward pass) consume memory proportional to $b \times s \times d$---the same on every GPU, unreduced by TP.

\paragraph{The Solution.}
\label{the-solution.}

Split the \emph{sequence dimension} for operations that don’t require cross-GPU communication (LayerNorm, Dropout, residual connections). Each GPU processes a $s/T$ slice of the sequence for these operations, then gathers the full sequence only where needed (attention, linear layers).

\begin{figure}[ht!]
\centering
\includegraphics[width=0.85\textwidth]{figures/fig_036_seq-parallel.png}
\caption{Sequence Parallelism reduces activation memory for LayerNorm/Dropout by splitting along the sequence dimension. Communication (AllGather/ReduceScatter) replaces the AllReduce used in standard TP---same total bytes transferred, but memory is saved.}
\label{fig:seq-parallel}
\end{figure}

\begin{intuitionbox}[SP Communication is ``Free'']
Standard TP uses AllReduce after each sub-layer, which is equivalent to ReduceScatter + AllGather. SP simply \emph{reorders} these primitives:

\begin{itemize}
  \item TP without SP: AllReduce (= ReduceScatter + AllGather) $\rightarrow$ same data on all GPUs $\rightarrow$ LayerNorm on full tensor (wasteful).
  \item TP with SP: ReduceScatter $\rightarrow$ each GPU has $1/T$ of sequence $\rightarrow$ LayerNorm on partial tensor $\rightarrow$ AllGather before next TP layer.
\end{itemize}

The total communication volume is identical! SP is purely a memory optimization with \textbf{zero additional communication cost}. It should always be enabled when using TP.
\end{intuitionbox}

\textbf{Memory savings from SP (70B model, TP=8, batch=4, seq=2048)}: 
\begin{equation}
\text{Activation savings} = (T-1) \times b \times s \times d \times n_\text{layers} \times 2\text{ bytes} = 7 \times 4 \times 2048 \times 8192 \times 80 \times 2 \approx \textbf{59 GB/GPU}
\end{equation}

\subsection{Pipeline Parallelism (PP)}
\label{subsec:pipeline-parallel}

Pipeline Parallelism splits the model \emph{vertically} by layers, assigning consecutive groups of layers to different devices (stages). Activations flow forward through stages; gradients flow backward.

\paragraph{The Bubble Problem.}
\label{the-bubble-problem.}

Naive pipeline execution creates ``bubbles''---idle time while a stage waits for input from the previous stage or gradients from the next:

\begin{figure}[ht!]
\centering
\includegraphics[width=0.95\textwidth]{figures/fig_037_pipeline-bubble.png}
\caption{Pipeline bubble comparison. Left: naive pipeline with one micro-batch has 75\% idle time. Right: GPipe with $M=4$ micro-batches reduces bubbles significantly. With $M \gg P$, bubble fraction approaches zero.}
\label{fig:pipeline-bubble}
\end{figure}

\paragraph{Bubble Fraction Formula.}
\label{bubble-fraction-formula.}

For $P$ pipeline stages and $M$ micro-batches per step: 
\begin{equation}
\text{Bubble fraction} = \frac{P - 1}{P + M - 1} \approx \frac{P-1}{M} \quad \text{(when } M \gg P\text{)}
\end{equation}

To keep bubble overhead $<$10\%, you need $M \geq 10 \cdot (P-1)$. For PP=4: at least 30 micro-batches.

\paragraph{Pipeline Schedules.}
\label{pipeline-schedules.}

\begin{table}[ht!]
\centering
\caption{Pipeline scheduling strategies}
\begin{tabular}{@{}lp{3.5cm}p{3.5cm}p{6cm}@{}}
\toprule
\textbf{Schedule} & \textbf{Bubble} & \textbf{Memory} & \textbf{Characteristics} \\
\midrule
GPipe & $\frac{P-1}{M+P-1}$ & $M \times$ activations & Simple; all-forward then all-backward~\cite{huang2019gpipe} \\
1F1B & $\frac{P-1}{M+P-1}$ & $P \times$ activations & Interleaved; steady-state memory bounded~\cite{narayanan2019pipedream} \\
Interleaved 1F1B & $\frac{P-1}{M \cdot V + P - 1}$ & $P \times$ activations & Virtual stages ($V$); further reduces bubble~\cite{narayanan2021efficient} \\
Zero-Bubble (ZB-H1) & $\approx 0$ & $P \times$ activations & Splits backward into B and W phases~\cite{qi2023zerobubble} \\
\bottomrule
\end{tabular}
\end{table}

\begin{intuitionbox}[1F1B: The Production Standard]
The \textbf{1F1B} (one-forward-one-backward) schedule~\cite{narayanan2019pipedream} is used in most production systems (Megatron-LM~\cite{narayanan2021efficient}, DeepSpeed~\cite{rajbhandari2020zero}):

\textbf{Warmup}: Forward passes fill the pipeline (P-1 micro-batches).

\textbf{Steady state}: Alternate one forward and one backward per time slot. This bounds peak activation memory to $P$ micro-batches (vs $M$ for GPipe).

\textbf{Cooldown}: Remaining backward passes drain the pipeline.

\textbf{Memory advantage}: GPipe must store activations for \emph{all} $M$ micro-batches simultaneously. 1F1B only stores $P$ sets of activations at steady state---critical when $M = 32$ but $P = 4$.
\end{intuitionbox}

\paragraph{Communication in PP.}
\label{communication-in-pp.}

Unlike TP (AllReduce), PP only requires \textbf{point-to-point} communication of activations between adjacent stages: 
\begin{equation}
\text{Data per transfer} = b_\text{micro} \times s \times d \times 2\text{ bytes (BF16)}
\end{equation}
 For micro-batch=4, seq=2048, $d$=8192: $4 \times 2048 \times 8192 \times 2 = 128$~MB per transfer. At InfiniBand 50~GB/s: 2.6~ms per transfer---small relative to compute per stage.

\paragraph{Load Balancing.}
\label{load-balancing.}

Not all layers have equal compute:

\begin{itemize}
  \item \textbf{Embedding layer}: Very cheap (lookup table).
  \item \textbf{Transformer blocks}: Uniform compute.
  \item \textbf{Final LM head}: Moderate (large matrix multiply for vocabulary projection).
\end{itemize}

Assign more transformer layers to middle stages and fewer to the first/last stages to balance compute.

\subsection{Fully Sharded Data Parallelism (FSDP / ZeRO-3)}
\label{fully-sharded-data-parallelism-fsdp-zero-3}

FSDP~\cite{zhao2023pytorch} (PyTorch) and ZeRO-3~\cite{rajbhandari2020zero} (DeepSpeed) address the memory duplication inherent in DDP: instead of every GPU holding a full copy of parameters, gradients, and optimizer states, each GPU owns only a $1/N$ slice and reconstructs the full tensor on-the-fly when needed.

\begin{figure}[ht!]
\centering
\includegraphics[width=0.95\textwidth]{figures/fig_038_fsdp.png}
\caption{FSDP shards all model state across GPUs. Each GPU owns $1/N$ of parameters, optimizer states, and gradients. Full parameters are reconstructed on-demand via AllGather before each layer’s computation.}
\label{fig:fsdp}
\end{figure}

\textbf{FSDP execution flow per layer:}

\begin{enumerate}
  \item \textbf{Forward}: AllGather parameters $\rightarrow$ compute $\rightarrow$ discard non-owned shards.
  \item \textbf{Backward}: AllGather parameters (again) $\rightarrow$ compute gradients $\rightarrow$ ReduceScatter gradients (each GPU gets its gradient shard) $\rightarrow$ discard non-owned parameter shards.
  \item \textbf{Optimizer step}: Each GPU updates only its owned shard using its gradient shard and optimizer states.
\end{enumerate}

\begin{table}[ht!]
\centering
\caption{Memory comparison: DDP vs FSDP/ZeRO stages (70B model, 8 GPUs). Baseline: BF16 params (140 GB) + BF16 grads (140 GB) + FP32 master+m+v (840 GB) = 1120 GB per GPU.}
\begin{tabular}{@{}lp{3.5cm}p{3.5cm}p{6cm}@{}}
\toprule
\textbf{Strategy} & \textbf{Sharded} & \textbf{Memory/GPU} & \textbf{Communication} \\
\midrule
DDP (no sharding) & Nothing & 1120 GB $\times$ & AllReduce (gradients only) \\
ZeRO-1 & Optimizer states & 385 GB $\times$ & AllReduce (gradients) \\
ZeRO-2 & Optimizer + gradients & 368 GB $\times$ & AllReduce (gradients) \\
ZeRO-3 / FSDP & Everything & \textbf{140 GB} \checkmark{} & AllGather + ReduceScatter (per layer) \\
\bottomrule
\end{tabular}
\end{table}

\begin{lstlisting}[style=pythonstyle]
from functools import partial
from torch.distributed.fsdp import FullyShardedDataParallel as FSDP
from torch.distributed.fsdp import ShardingStrategy, MixedPrecision, BackwardPrefetch
from torch.distributed.fsdp.wrap import transformer_auto_wrap_policy
from transformers.models.llama.modeling_llama import LlamaDecoderLayer

# Wrap model with FSDP
auto_wrap = partial(transformer_auto_wrap_policy,
                    transformer_layer_cls={LlamaDecoderLayer})
mp_policy = MixedPrecision(
    param_dtype=torch.bfloat16,
    reduce_dtype=torch.bfloat16,
    buffer_dtype=torch.bfloat16,
)

model = FSDP(
    model,
    sharding_strategy=ShardingStrategy.FULL_SHARD,  # ZeRO-3
    mixed_precision=mp_policy,
    auto_wrap_policy=auto_wrap,  # Wrap each transformer layer
    use_orig_params=True,        # Required for torch.compile compatibility
    limit_all_gathers=True,      # Bound peak memory (1 AllGather in flight at a time)
    forward_prefetch=True,       # Prefetch next layer's params during current layer
    backward_prefetch=BackwardPrefetch.BACKWARD_PRE,  # Prefetch during backward
)
\end{lstlisting}

\begin{warningbox}[FSDP Communication Volume]
FSDP communicates \textbf{3$\times$ more data} than DDP per step:

\begin{itemize}
  \item DDP: 1 AllReduce of gradients = $2M$ bytes total across ring (where $M$ = model size in bytes).
  \item FSDP: 2 AllGather (forward + backward) + 1 ReduceScatter = $3M$ bytes.
\end{itemize}

This is the memory--communication trade-off. FSDP is worthwhile when: (a) model doesn’t fit in GPU memory with DDP, or (b) communication is well-overlapped with compute (modern frameworks achieve 70--90\% overlap).
\end{warningbox}

\subsection{3D Parallelism: Combining Strategies}
\label{subsec:3d-parallel}

Production systems at scale (70B+) combine TP, PP, and DP/FSDP simultaneously:

\begin{figure}[ht!]
\centering
\includegraphics[width=0.95\textwidth]{figures/fig_039_3d-parallel.png}
\caption{3D parallelism layout for 16 GPUs: TP=4 (within each box, using NVLink), PP=2 (orange arrows, stages), DP=2 (red arrows, gradient sync). Each dimension exploits a different level of the communication hierarchy.}
\label{fig:3d-parallel}
\end{figure}

\begin{keybox}[Production Recipe: 70B on 64 A100-80GB (8 nodes)]
\textbf{Intra-node} (NVLink 600GB/s): TP=8 for generation, FSDP within node for training.\\

\textbf{Inter-node} (InfiniBand 400Gb/s): FSDP across nodes (8-way data parallel).\\

\textbf{Result}: Each GPU holds $\sim$70GB. Policy weights gathered per-layer during forward/backward.\\

\textbf{Pipeline Parallel}: Only if model exceeds 100B+ and won’t fit with TP+ZeRO. Adds complexity (bubble overhead 10--20\%) and scheduling headaches.

\textbf{Decision flowchart}:

\begin{enumerate}
  \item Does the model fit on 1 GPU? $\rightarrow$ Use DDP.
  \item Does it fit on 1 node with FSDP? $\rightarrow$ Use FSDP (ZeRO-3).
  \item Does it fit on 1 node with TP+FSDP? $\rightarrow$ Use TP (intra-node) + FSDP (inter-node).
  \item Still doesn’t fit? $\rightarrow$ Add PP across nodes. This is the last resort.
\end{enumerate}
\end{keybox}

\begin{table}[ht!]
\centering
\caption{Parallelism strategy comparison summary}
{\footnotesize
\begin{tabular}{@{}llllll@{}}
\toprule
\textbf{Strategy} & \textbf{Splits} & \textbf{Communication} & \textbf{Scaling Limit} & \textbf{Overhead} & \textbf{When to Use} \\
\midrule
DP/DDP & Batch & AllReduce (grads) & $\sim$64 GPUs & 5--10\% & Model fits on 1 GPU \\
FSDP & Params+Opt+Grad & AllGather+RS & 100s of GPUs & 10--20\% & Default for $>$13B \\
TP & Weight matrices & AllReduce (2/layer) & 8 GPUs (1 node) & 12--18\% & Large model inference+train \\
SP & Activations (seq) & Reuses TP comms & Same as TP & $\approx$0\% extra & Always with TP \\
PP & Layers (stages) & Point-to-point & $\sim$16 stages & 15--30\% & 100B+ models only \\
\bottomrule
\end{tabular}
}
\end{table}

\section{The Generation Bottleneck: Quantitative Analysis}
\label{the-generation-bottleneck-quantitative-analysis}

\begin{intuitionbox}[Roofline Analysis: Why Generation is Memory-Bound]
\textbf{A100 specs}: 312 TFLOPS (BF16 tensor cores), 2 TB/s HBM bandwidth.

\textbf{Roofline crossover}: $312\text{T} / 2\text{T} = 156$ FLOP/byte. Operations below 156 FLOP/byte are \emph{memory-bound}.

\textbf{Autoregressive generation}: For each token, read all weights (140GB for 70B) and do $2 \times 70\text{B} = 140\text{G}$ FLOPs per token (at batch=1).

\textbf{Arithmetic intensity}: $140\text{G FLOP} / 140\text{GB} = 1$ FLOP/byte. That’s $156\times$ below the roofline!

\textbf{Utilization}: $1/156 = 0.6\%$ of peak FLOPS utilized. The GPU is 99.4\% idle, waiting for memory reads.

\textbf{Token rate}: $2\text{TB/s} / 140\text{GB} = 14.3$ tokens/second (single stream, batch=1).

\textbf{For 512 tokens}: $512 / 14.3 = 35.8$ seconds per response (batch=1, TP=1).

\textbf{Batching helps}: Batch=64 with TP=4 $\rightarrow$ reads weights once, generates 64 tokens in parallel. Arithmetic intensity: $64 \times 1 = 64$ FLOP/byte. Better, but still below roofline!
\end{intuitionbox}

\begin{table}[ht!]
\centering
\caption{Generation throughput for 70B model (512 tokens, various configurations)}
{\footnotesize
\begin{tabular}{@{}lllll@{}}
\toprule
\textbf{Config} & \textbf{Batch} & \textbf{Time/batch} & \textbf{Tok/s/GPU} & \textbf{Notes} \\
\midrule
TP=1, batch=1 & 1 & 36s & 14 & Baseline, worst case \\
TP=4, batch=1 & 1 & 9s & 57 & Linear TP scaling for gen \\
TP=4, batch=32 & 32 & 15s & 1092 & Near-optimal batching \\
TP=4, batch=128, vLLM & 128 & 45s & 1456 & Continuous batching \\
TP=4, batch=128, INT8 & 128 & 25s & 2621 & 2$\times$ bandwidth savings \\
\bottomrule
\end{tabular}
}
\end{table}

\textbf{Optimization stack} (cumulative speedup):

\begin{enumerate}
  \item \textbf{vLLM + PagedAttention}~\cite{kwon2023efficient} (2--4$\times$): Eliminates KV cache fragmentation, enables larger batches
  \item \textbf{Continuous batching}~\cite{yu2022orca} (1.5--2$\times$): Don’t wait for longest sequence; start new ones as others finish
  \item \textbf{Speculative decoding}~\cite{leviathan2023fast} (2--3$\times$): Small draft model proposes 5 tokens, large model verifies in one forward pass. Accept 3--4 on average.
  \item \textbf{INT8/FP8 weights for gen} (2$\times$): Halve bandwidth needs. Quality loss is minimal since we’re sampling (not computing exact logits for training)
  \item \textbf{CUDA graphs} (1.1--1.3$\times$): Eliminate kernel launch overhead for fixed-shape operations
  \item \textbf{Prefix caching} (1.5$\times$ for shared-prefix prompts): Don’t recompute system prompt KV cache
\end{enumerate}

\begin{lstlisting}[style=pythonstyle]
# Production vLLM generation setup
from vllm import LLM, SamplingParams

engine = LLM(
    model="./policy_checkpoint",
    tensor_parallel_size=4,           # TP=4 per instance
    gpu_memory_utilization=0.92,      # Leave headroom for KV cache
    max_num_batched_tokens=16384,     # Max tokens in flight
    max_num_seqs=256,                 # Max concurrent sequences
    dtype="bfloat16",
    enable_prefix_caching=True,       # Cache system prompt KV
    speculative_model="./draft_1B",   # Speculative decoding
    num_speculative_tokens=5,
    block_size=16,                    # PagedAttention block size
    swap_space=4,                     # GB swap space for preemption
)

# Generate responses for RLHF batch
sampling_params = SamplingParams(
    temperature=0.7, top_p=0.9, max_tokens=512,
    logprobs=1,  # Need log-probs for PPO ratio calculation
)
outputs = engine.generate(prompts, sampling_params)
# Extract: responses, log_probs for each token (needed for PPO/GRPO)
\end{lstlisting}

\section{Decoupled Architecture: Production Design}
\label{decoupled-architecture-production-design}

Production RLHF systems such as DeepSpeed-Chat~\cite{yao2023deepspeedchat} and OpenRLHF~\cite{hu2024openrlhf} use a \textbf{decoupled architecture} that separates generation, scoring, and training into independently scalable clusters.

\begin{figure}[ht!]
\centering
\includegraphics[width=0.85\textwidth]{figures/fig_040_fig40.png}
\caption{Decoupled RLHF architecture. Each cluster optimized for its workload. Scored rollouts accumulate in the experience buffer before being consumed by training.}
\end{figure}

\begin{keybox}[Why Decouple?]
\textbf{Generation} is memory-bandwidth bound (need fast HBM, waste compute).\\

\textbf{Training} is compute-bound (need tensor cores, waste bandwidth during backprop).\\

\textbf{Same hardware can’t optimize both}: If you put everything together, you either waste compute during generation or waste bandwidth during training. Decoupling lets each cluster use optimal hardware/config.

\textbf{Practical benefits}:

\begin{itemize}
  \item Scale generation and training independently
  \item Generation cluster is stateless $\rightarrow$ trivial fault tolerance
  \item Can overlap gen(step $N+1$) with training(step $N$) $\rightarrow$ 30--40\% speedup
  \item Different quantization: INT8 for generation (bandwidth), BF16 for training (precision)
\end{itemize}
\end{keybox}

\section{Weight Synchronization Strategies}
\label{weight-synchronization-strategies}

{\small
\begin{tabular}{@{}llll@{}}
\toprule
\textbf{Strategy} & \textbf{Staleness} & \textbf{Bandwidth} & \textbf{Quality Impact} \\
\midrule
Synchronous (every step) & 0 steps & 140 GB/step & Perfect but too slow \\
Periodic (every 50 steps) & 25 avg & 2.8 GB/step amortized & $<$2\% quality loss \\
Delta compression (INT8) & 25 avg & 0.4 GB/step & $<$3\% quality loss \\
Async streaming & 5--10 steps & 14 GB/step (background) & $<$1\% quality loss \\
\bottomrule
\end{tabular}
}

\begin{intuitionbox}[Why Staleness is OK for PPO/GRPO]
PPO’s clipped objective was \emph{designed} for off-policy data! The clip $[1-\epsilon, 1+\epsilon]$ bounds the impact of stale data. With 10--50 steps of staleness:

\begin{itemize}
  \item Policy changes $\sim$0.1--1\% per step (with proper LR)
  \item Over 50 steps: $\sim$5\% policy drift
  \item PPO clip handles up to 20\% drift by design
  \item Empirically: quality loss $<$2\% for 50-step staleness
\end{itemize}

\textbf{Bandwidth math}: 70B BF16 = 140GB. InfiniBand 400Gb/s = 50GB/s $\rightarrow$ full sync in 2.8s. With delta compression: $<$0.5s. Async = free (runs in background).
\end{intuitionbox}

\section{Memory Optimization Techniques}
\label{memory-optimization-techniques}

\begin{tabular}{@{}lp{5cm}p{8cm}@{}}
\toprule
\textbf{ZeRO Stage} & \textbf{What Gets Sharded} & \textbf{Memory/GPU (70B, 8 GPUs)} \\
\midrule
None (Data Parallel) & Nothing (full replica) & 560GB per GPU (impossible) \\
ZeRO-1 & Optimizer states only & 175GB \\
ZeRO-2 & Optimizer states + Gradients & 105GB \\
ZeRO-3 (FSDP) & Optimizer + Gradients + Parameters & \textbf{70GB (fits in A100-80GB!)} \\
\bottomrule
\end{tabular}

\textbf{Additional techniques}:

\begin{itemize}
  \item \textbf{Gradient checkpointing}~\cite{chen2016training}: Don’t store all activations; recompute during backward pass. Saves $\sim$60\% activation memory, costs $\sim$33\% extra compute. Selective: only checkpoint attention layers (memory-heavy), keep FFN activations (compute-heavy to recompute).
  \item \textbf{Mixed precision}~\cite{micikevicius2018mixed}: Forward in BF16 (2 bytes/param), optimizer states in FP32 (4 bytes each for m,v). Master weights in FP32 for accumulation.
  \item \textbf{CPU offloading} (ZeRO-Infinity~\cite{rajbhandari2021zeroinfinity}): Move optimizer states to CPU RAM. 50\% memory savings but 2--3$\times$ slower (PCIe 64GB/s bottleneck).
  \item \textbf{Activation offloading}: Move activations to CPU during forward, bring back for backward. Only when memory is truly critical.
  \item \textbf{Flash Attention}~\cite{dao2022flashattention, dao2023flashattention2}: O($n$) memory instead of O($n^2$) for attention. 2--4$\times$ faster + massive memory savings for long sequences.
\end{itemize}

\subsection{Flash Attention’s Impact on RLHF}
\label{flash-attentions-impact-on-rlhf}

\begin{intuitionbox}[Why Flash Attention Matters for RLHF]
RLHF involves generating long sequences (rollouts) and then training on them. Without Flash Attention:

\begin{itemize}
  \item A 4K-token sequence with 32 heads requires $\sim$4 GB just for attention matrices
  \item This severely limits batch size during PPO/GRPO training
  \item Gradient checkpointing of attention activations is expensive
\end{itemize}

With Flash Attention:

\begin{itemize}
  \item Attention memory is $O(n)$ -- dominated by $Q, K, V, O$ tensors
  \item Longer rollouts (8K--32K tokens) become feasible with the same GPU memory
  \item Backward pass recomputes attention tiles from $Q, K, V$ (no stored $n^2$ matrix)
  \item This is the key enabler for long-context RLHF (e.g., reasoning models)
\end{itemize}
\end{intuitionbox}

\begin{warningbox}[Flash Attention and Gradient Checkpointing]
Flash Attention’s backward pass recomputes the attention tiles on-the-fly from $Q, K, V$ (which are stored). This means Flash Attention \emph{already implements} a form of activation recomputation for the $O(n^2)$ attention matrix. You do not need to additionally checkpoint the attention layer -- doing so would recompute $Q, K, V$ unnecessarily.
\end{warningbox}

\begin{lstlisting}[style=pythonstyle]
# DeepSpeed ZeRO-3 configuration for 70B RLHF training
ds_config = {
    "bf16": {"enabled": True},
    "zero_optimization": {
        "stage": 3,
        "overlap_comm": True,                    # Overlap communication with compute
        "contiguous_gradients": True,            # Better memory layout
        "reduce_scatter": True,                  # More efficient than allreduce
        "reduce_bucket_size": 5e7,               # 50M params per bucket
        "prefetch_bucket_size": 5e7,             # Prefetch next bucket
        "param_persistence_threshold": 1e5,      # Keep small params on all GPUs
        "offload_optimizer": {"device": "cpu", "pin_memory": True},  # CPU offload
        "sub_group_size": 1e9,                   # Reduce fragmentation
    },
    "gradient_accumulation_steps": 4,
    "gradient_clipping": 1.0,
    "train_micro_batch_size_per_gpu": 2,
    "wall_clock_breakdown": True,
}
\end{lstlisting}

\section{Fault Tolerance at Scale}
\label{fault-tolerance-at-scale}

\begin{warningbox}[Hardware Failure Reality]
\textbf{Individual GPU MTBF}: $\sim$10,000 hours.\\
\textbf{512-GPU cluster MTBF}: $10000/512 \approx 20$ hours. But with software/network: \textbf{4--8 hours realistically}.\\
\textbf{Multi-day training run}: Will see 5--15 failures. Without fault tolerance, one failure kills everything.
\end{warningbox}

\textbf{Production fault tolerance stack}:

\begin{enumerate}
  \item \textbf{Detection}: NCCL timeout (60s), GPU heartbeat (10s), NVML health monitoring, ECC error counting.
  \item \textbf{Checkpointing}: Async every 50--100 steps. Non-blocking (background thread). Save: model weights, optimizer states (Adam m/v), scheduler state, RNG states, KL coefficient, replay buffer. Keep last 3 checkpoints. Time: $\sim$30s for 70B (parallel write to NVMe).
  \item \textbf{Recovery}: (a) Generation cluster = stateless, just restart and load latest weights. (b) Training cluster: load checkpoint, rebuild NCCL process group excluding failed node, redistribute FSDP shards, resume from last checkpoint.
  \item \textbf{Elastic training}: Torch Elastic / Kubernetes auto-scaling. Replace failed node within minutes. Training continues with $N-1$ GPUs temporarily.
  \item \textbf{Prevention}: GPU health pre-screening (run GEMM stress test before starting). Hot spares on standby. Redundant network paths (dual-rail InfiniBand).
\end{enumerate}

\section{End-to-End Latency Breakdown}
\label{end-to-end-latency-breakdown}

\begin{figure}[ht!]
\centering
\includegraphics[width=0.85\textwidth]{figures/fig_041_fig41.png}
\caption{Without overlap (monolithic). With decoupled: gen overlaps with training, effective 1.4$\times$ speedup.}
\end{figure}

\begin{tabular}{@{}lp{3.5cm}p{3.5cm}p{6cm}@{}}
\toprule
\textbf{Phase} & \textbf{Time (70B)} & \textbf{Bound By} & \textbf{Optimization} \\
\midrule
Generation (128$\times$512 tok) & 30--45s & Memory bandwidth & vLLM, spec decoding, INT8 \\
Reward scoring & 5--8s & Compute (batch forward) & INT8 RM, batch=128 \\
Reference log-probs & 4--6s & Compute (batch forward) & INT8 ref, or LoRA (free) \\
PPO update (4 epochs) & 8--12s & Compute (backprop) & FSDP, Flash Attention \\
Weight sync & 0--3s & Network (async) & Delta compression, async \\
\textbf{Total (monolithic)} & \textbf{50--75s} &  &  \\
\textbf{Total (decoupled, overlapped)} & \textbf{35--50s} &  & Gen overlaps with prev training \\
\bottomrule
\end{tabular}

\section{Monitoring and Observability}
\label{monitoring-and-observability}

\begin{keybox}[Key Metrics to Track During RLHF Training]
\textbf{Quality metrics} (log every 10 steps):

\begin{itemize}
  \item Mean reward (should increase then plateau)
  \item KL divergence from reference (should stay 3--10)
  \item Response length distribution (watch for length hacking)
  \item Entropy (should decrease slowly, not collapse)
\end{itemize}

\textbf{System metrics} (log every step):

\begin{itemize}
  \item GPU utilization (target: $>$80\% during training, $>$60\% during gen)
  \item Memory watermark per GPU (catch OOM before it happens)
  \item Generation throughput (tokens/sec, should be stable)
  \item Gradient norm (spikes = instability incoming)
  \item NCCL communication time (detect network degradation)
\end{itemize}
\end{keybox}

\section{Network Topology and Communication Patterns}
\label{sec:network-topology}

Efficient distributed training requires understanding the hierarchical communication fabric that connects GPUs. Modern clusters use a two-tier architecture: ultra-fast intra-node links and slower but scalable inter-node networks.

\subsection{Intra-Node: NVLink and NVSwitch}
\label{intra-node-nvlink-and-nvswitch}

\begin{table}[ht!]
\centering
\caption{NVLink generations and their impact on LLM training}
\begin{tabular}{@{}lp{3cm}p{3.2cm}p{3.2cm}p{3.5cm}@{}}
\toprule
\textbf{Generation} & \textbf{BW per link} & \textbf{Links/GPU} & \textbf{Total BW} & \textbf{Platform} \\
\midrule
NVLink 3.0 & 50 GB/s & 12 & 600 GB/s & A100 (DGX A100) \\
NVLink 4.0 & 50 GB/s & 18 & 900 GB/s & H100 (DGX H100) \\
NVLink 5.0 & 100 GB/s & 18 & 1800 GB/s & B200 (DGX B200) \\
\bottomrule
\end{tabular}
\end{table}

Within a single node (typically 8 GPUs), \textbf{NVSwitch} provides full-bisection bandwidth between all GPU pairs. This means any GPU can communicate with any other at full NVLink speed simultaneously---critical for Tensor Parallelism where every layer requires an AllReduce across all 8 GPUs.

\begin{keybox}[NVSwitch vs PCIe Topology]
\textbf{With NVSwitch} (DGX/HGX): All 8 GPUs connected all-to-all at 600--1800~GB/s. AllReduce for TP takes $\sim$0.2ms per layer.

\textbf{Without NVSwitch} (PCIe-only servers): GPUs communicate through CPU PCIe root complex at 32--64~GB/s. TP across 8 GPUs becomes 10--30$\times$ slower. \textbf{Never use TP$>$2 on PCIe-only systems.}
\end{keybox}

\subsection{Inter-Node: InfiniBand and RoCE}
\label{inter-node-infiniband-and-roce}

For FSDP/ZeRO-3 AllGather and ReduceScatter operations across nodes, the inter-node network dominates.

\begin{table}[ht!]
\centering
\caption{Inter-node networking options for LLM training clusters}
{\small
\begin{tabular}{@{}llll@{}}
\toprule
\textbf{Technology} & \textbf{Bandwidth} & \textbf{Latency} & \textbf{Notes} \\
\midrule
InfiniBand NDR & 400 Gb/s (50 GB/s) & 1--2 $\mu$s & Gold standard, RDMA, lossless \\
InfiniBand NDR (dual-rail) & 800 Gb/s (100 GB/s) & 1--2 $\mu$s & Used in H100 clusters \\
RoCE v2 & 100--400 Gb/s & 2--5 $\mu$s & Cheaper, needs PFC/ECN tuning \\
Ethernet (TCP) & 100--400 Gb/s & 10--50 $\mu$s & Not suitable for $>$16 GPU training \\
\bottomrule
\end{tabular}
}
\end{table}

\subsection{Communication Primitives and Their Costs}
\label{communication-primitives-and-their-costs}

Understanding when each collective is used helps diagnose bottlenecks:

\begin{table}[ht!]
\centering
\caption{NCCL collective operations in distributed LLM training}
\begin{tabular}{@{}lp{3.5cm}p{3.5cm}p{6cm}@{}}
\toprule
\textbf{Collective} & \textbf{Data Moved} & \textbf{Used By} & \textbf{When} \\
\midrule
AllReduce & $2 \cdot \frac{N-1}{N} \cdot M$ & TP, DP & Sum gradients or activations across GPUs \\
AllGather & $\frac{N-1}{N} \cdot M$ & FSDP forward & Reconstruct full parameter tensor before matmul \\
ReduceScatter & $\frac{N-1}{N} \cdot M$ & FSDP backward & Distribute gradient shards after backprop \\
Broadcast & $M$ & PP & Send activations to next pipeline stage \\
Send/Recv & $M$ & PP & Point-to-point between adjacent stages \\
\bottomrule
\end{tabular}
\end{table}

where $M$ is the message size (bytes) and $N$ is the number of participants.

\newpage
\begin{intuitionbox}[Communication-Computation Overlap]
Modern frameworks (FSDP, DeepSpeed) aggressively overlap communication with computation:

\textbf{Forward pass}: While layer $i$ computes, AllGather prefetches parameters for layer $i+1$. After layer $i$ finishes, its parameters are immediately discarded (``free-after-forward'').

\textbf{Backward pass}: While layer $i$ computes gradients, ReduceScatter sends layer $i+1$’s gradients. This overlap hides 70--90\% of communication latency when properly tuned.

\textbf{Tuning knobs}: \texttt{prefetch\_factor} (how many layers ahead to prefetch), \texttt{reduce\_bucket\_size} (granularity of gradient reduction), \texttt{backward\_prefetch} (``pre'' vs ``post'' backward prefetch strategy).
\end{intuitionbox}

\subsection{Network Topology Design}
\label{network-topology-design}

Production clusters use \textbf{fat-tree} or \textbf{rail-optimized} topologies:

\begin{itemize}
  \item \textbf{Fat-tree}: Full bisection bandwidth at every level. Any node can communicate with any other at full speed. Expensive (many switches) but maximally flexible.
  \item \textbf{Rail-optimized}: GPU $i$ on every node connects to the same leaf switch (``rail $i$''). AllReduce within a rail is cheap; cross-rail traffic is expensive. Used by Meta’s RSC and Google’s TPU pods.
  \item \textbf{3D torus / Dragonfly}: Used in HPC clusters (Frontier, Aurora). Topology-aware job placement is critical.
\end{itemize}

\begin{warningbox}[Job Placement Matters]
On a 512-GPU cluster, random node assignment can cause 2--3$\times$ slowdown due to network congestion. \textbf{Always request contiguous node blocks.} Production schedulers (Slurm, Kubernetes) should enforce locality: all nodes in a training job should be on the same leaf switch or within one hop of each other.
\end{warningbox}

\section{Training Throughput and Model FLOPs Utilization}
\label{sec:mfu}

\subsection{Measuring Training Efficiency: MFU}
\label{measuring-training-efficiency-mfu}

\textbf{Model FLOPs Utilization (MFU)}~\cite{chowdhery2022palm} is the standard metric for training efficiency:

\begin{equation}
\text{MFU} = \frac{\text{Observed throughput (tokens/sec)} \times \text{FLOPs per token}}{\text{Peak hardware FLOPS}}
\end{equation}

For a transformer with $P$ parameters, $s$ sequence length, and $b$ batch size: 
\begin{equation}
\text{FLOPs per token} \approx 6P + 12 \cdot n_\text{layers} \cdot d_\text{model} \cdot s
\end{equation}

The factor of 6 comes from: 2 (multiply-add) $\times$ 3 (forward + backward, where backward $\approx 2\times$ forward). The second term accounts for attention’s $O(s^2)$ cost.

\begin{table}[ht!]
\centering
\caption{MFU benchmarks across scales and hardware}
\begin{tabular}{@{}lp{3cm}p{3.2cm}p{3.2cm}p{3.5cm}@{}}
\toprule
\textbf{Model} & \textbf{Hardware} & \textbf{MFU} & \textbf{Tokens/sec/GPU} & \textbf{Configuration} \\
\midrule
LLaMA-7B & 8$\times$A100 & 57\% & 3,200 & FSDP, FlashAttn, BF16 \\
LLaMA-13B & 16$\times$A100 & 52\% & 1,750 & FSDP, FlashAttn, BF16 \\
LLaMA-70B & 64$\times$A100 & 45\% & 380 & FSDP+TP=8, FlashAttn \\
GPT-4 (est.) & 10,000+ H100 & 40--50\% & --- & 3D parallelism \\
PaLM-540B & 6144 TPUv4 & 46\% & --- & DP+TP+PP \\
\bottomrule
\end{tabular}
\end{table}

\begin{intuitionbox}[Why MFU Decreases with Scale]
Larger models require more parallelism, which introduces:

\begin{enumerate}
  \item \textbf{Communication overhead}: AllGather/ReduceScatter for FSDP ($\sim$10--15\% at 64 GPUs)
  \item \textbf{Pipeline bubbles}: PP introduces idle time at start/end of micro-batches ($\sim$15--25\% with PP=4)
  \item \textbf{Memory for auxiliary models}: Reference/RM take GPU memory that could hold larger batches
  \item \textbf{Load imbalance}: Not all layers have equal compute (embeddings vs transformer blocks)
\end{enumerate}

\textbf{Rule of thumb}: Target MFU $>$ 40\% for training. If below 30\%, diagnose with profiling.
\end{intuitionbox}

\subsection{Compute-Optimal Batch Sizing}
\label{compute-optimal-batch-sizing}

The effective batch size interacts with hardware utilization in non-obvious ways:

\begin{equation}
\text{Effective batch size} = \text{micro\_batch} \times \text{grad\_accum} \times \text{DP degree}
\end{equation}

\begin{itemize}
  \item \textbf{Too small}: GPU underutilized (low arithmetic intensity), communication dominates.
  \item \textbf{Too large}: Diminishing learning per token (critical batch size exceeded), wastes compute.
  \item \textbf{Sweet spot}: The \emph{critical batch size} $B_\text{crit}$ where gradient noise equals gradient signal. For LLMs, $B_\text{crit} \sim 1$--$4$M tokens~\cite{mccandlish2018empirical}.
\end{itemize}

For RLHF specifically, the batch contains \emph{rollouts} (not just tokens): 
\begin{equation}
\text{RLHF batch} = N_\text{prompts} \times K_\text{generations} \times L_\text{avg response length}
\end{equation}

Typical production values: $N=128$ prompts, $K=1$--$4$ generations, $L=256$--$512$ tokens $\rightarrow$ 32K--256K tokens per step.

\subsection{Profiling and Bottleneck Diagnosis}
\label{profiling-and-bottleneck-diagnosis}

Key profiling tools and what they reveal:

\begin{tabular}{@{}lp{5cm}p{8cm}@{}}
\toprule
\textbf{Tool} & \textbf{Captures} & \textbf{Best For} \\
\midrule
\texttt{torch.profiler} & Kernel timing, memory & Finding slow ops, memory leaks \\
NVIDIA Nsight Systems & Full GPU timeline & Visualizing overlap, gaps between kernels \\
\texttt{nccl\_debug=INFO} & Collective sizes/times & Diagnosing communication bottlenecks \\
\texttt{torch.cuda.memory\_stats} & Allocation patterns & Finding fragmentation, peak usage \\
DeepSpeed Flops Profiler & Per-layer FLOPs & Identifying load imbalance \\
\texttt{py-spy} / \texttt{scalene} & CPU profiling & Data loading, tokenization bottlenecks \\
\bottomrule
\end{tabular}

\begin{examplebox}[Diagnosing Low MFU: A Checklist]
\begin{enumerate}
  \item \textbf{GPU utilization $<$ 80\%?} $\rightarrow$ Data loading bottleneck (check CPU, I/O).
  \item \textbf{Large gaps between kernels?} $\rightarrow$ Python overhead, synchronization points. Use CUDA graphs.
  \item \textbf{Communication $>$ 20\% of step time?} $\rightarrow$ Reduce TP degree, increase batch size, check network health.
  \item \textbf{Memory at 99\%?} $\rightarrow$ Cannot increase batch. Try gradient checkpointing, offloading.
  \item \textbf{OOM during generation?} $\rightarrow$ KV cache too large. Reduce max\_seq\_len or batch size for gen.
\end{enumerate}
\end{examplebox}

\section{Cost Analysis and Cloud Deployment}
\label{cost-analysis-and-cloud-deployment}

Understanding the economics of RLHF training is essential for planning.

\subsection{Hardware Cost Comparison}
\label{hardware-cost-comparison}

\begin{table}[ht!]
\centering
\caption{Approximate cloud GPU costs for RLHF training (2024--2025 pricing)}
\begin{tabular}{@{}lp{3cm}p{3.2cm}p{3.2cm}p{3.5cm}@{}}
\toprule
\textbf{GPU} & \textbf{On-Demand/hr} & \textbf{Spot/hr} & \textbf{Memory} & \textbf{Use Case} \\
\midrule
A100 80GB & \$2.50--3.50 & \$1.00--1.50 & 80 GB HBM2e & Budget training, gen cluster \\
H100 80GB & \$4.00--6.00 & \$2.00--3.00 & 80 GB HBM3 & Production training \\
H200 141GB & \$6.00--8.00 & --- & 141 GB HBM3e & Large context, fewer-GPU configs \\
MI300X 192GB & \$3.50--5.00 & \$1.50--2.50 & 192 GB HBM3 & Cost-effective alternative \\
\bottomrule
\end{tabular}
\end{table}

\subsection{RLHF Training Cost Estimation}
\label{rlhf-training-cost-estimation}

\begin{equation}
\text{Cost} = \frac{N_\text{steps} \times T_\text{step}}{3600} \times N_\text{GPUs} \times C_\text{GPU/hr}
\end{equation}

\begin{examplebox}[Cost Example: 70B Model RLHF (10K steps)]
\begin{tabular}{@{}lp{11cm}@{}}
\toprule
Steps & 10,000 \\
\midrule
Time per step (decoupled) & 45 seconds \\
Total training time & $10000 \times 45 / 3600 = 125$ hours \\
GPUs (generation + training) & 64 A100-80GB \\
Cost per GPU-hour (spot) & \$1.20 \\
\textbf{Total cost} & $125 \times 64 \times \$1.20 =$ \textbf{\$9,600} \\
\bottomrule
\end{tabular}

\textbf{Breakdown by phase}:

\begin{itemize}
  \item Generation cluster (32 GPUs): \$4,800 (60\% of time)
  \item Training cluster (32 GPUs): \$4,800 (could overlap $\rightarrow$ \$3,400 effective)
  \item Scoring (shared with gen GPUs): included above
\end{itemize}

\textbf{With overlap}: Effective cost $\approx$ \textbf{\$7,500} for full RLHF alignment of a 70B model.
\end{examplebox}

\subsection{Cost Optimization Strategies}
\label{cost-optimization-strategies}

\begin{itemize}
  \item \textbf{Spot/preemptible instances}: 50--70\% savings. Requires robust checkpointing (save every 5 minutes).
  \item \textbf{Right-sizing}: Don’t use H100 for generation (memory-bound); A100 achieves similar tokens/\$ for inference.
  \item \textbf{Quantized inference}: INT8/FP8 for generation and scoring halves GPU count for those clusters.
  \item \textbf{Progressive training}: Start with 8B proxy model for reward engineering/debugging ($\sim$\$200), then scale to 70B.
  \item \textbf{LoRA for reference-free}: Eliminates reference model entirely (50\% memory reduction).
  \item \textbf{Shorter sequences first}: Curriculum from 256$\rightarrow$512$\rightarrow$1024 token generations saves 40\% compute.
\end{itemize}

\section{Distributed Checkpointing}
\label{sec:checkpointing}

At scale, naive checkpointing becomes a bottleneck. A 70B model with optimizer state requires saving $\sim$840~GB per checkpoint (FP32 master weights + Adam m + v).

\subsection{Checkpointing Strategies}
\label{checkpointing-strategies}

\begin{table}[ht!]
\centering
\caption{Checkpointing approaches for large-scale RLHF}
\begin{tabular}{@{}lp{3.5cm}p{3.5cm}p{6cm}@{}}
\toprule
\textbf{Strategy} & \textbf{Save Time (70B)} & \textbf{Storage/ckpt} & \textbf{Characteristics} \\
\midrule
Synchronous (all ranks) & 30--60s (blocking) & 420 GB & Simple, stalls training \\
Async (background copy) & $<$1s (non-blocking) & 420 GB & Overlaps with next step \\
Incremental (delta) & $<$1s & 5--20 GB & Only save changed params \\
Sharded (FSDP native) & 5--10s & 420 GB sharded & Each rank saves its shard \\
\bottomrule
\end{tabular}
\end{table}

\subsection{Production Checkpointing with torch.distributed.checkpoint}
\label{production-checkpointing-with-torch.distributed.checkpoint}

\begin{lstlisting}[style=pythonstyle]
import torch.distributed.checkpoint as dcp
from torch.distributed.checkpoint.state_dict import get_state_dict, StateDictOptions

# Save: each rank writes its shard in parallel
state_dict = {"model": get_state_dict(model, options=StateDictOptions(full_state_dict=False))}
dcp.save(
    state_dict=state_dict,
    storage_writer=dcp.FileSystemWriter("/mnt/checkpoints/step_5000"),
    planner=dcp.DefaultSavePlanner(),  # Handles FSDP sharding automatically
)

# Async save: non-blocking, runs in background thread
future = dcp.async_save(
    state_dict=state_dict,
    storage_writer=dcp.FileSystemWriter("/mnt/checkpoints/step_5000"),
)
# Training continues immediately; future.result() blocks only if needed
\end{lstlisting}

\begin{keybox}[Checkpoint Hygiene for RLHF]
RLHF checkpoints must capture \emph{more} than standard pre-training:

\begin{itemize}
  \item Policy model weights + optimizer states (standard)
  \item KL coefficient ($\beta$) and its schedule state
  \item Replay buffer contents (for off-policy corrections)
  \item RNG states for all GPUs (reproducibility)
  \item Prompt iterator position (avoid re-processing prompts)
  \item Reward model version tag (for auditability)
  \item Wandb/metrics run ID (for continuous logging)
\end{itemize}
\end{keybox}

\section{Hardware Selection Guide}
\label{sec:hardware-guide}

Choosing the right hardware depends on model size, budget, and training phase.

\begin{table}[ht!]
\centering
\caption{Hardware recommendations by model scale and training phase}
\begin{tabular}{@{}lp{3.5cm}p{3.5cm}p{6cm}@{}}
\toprule
\textbf{Model Size} & \textbf{Training Phase} & \textbf{Recommended} & \textbf{Configuration} \\
\midrule
$\leq$7B & SFT + RLHF & 1--2$\times$ A100 & Single node, no parallelism needed \\
7--13B & SFT + RLHF & 4--8$\times$ A100 & FSDP, optional TP=2 for gen \\
13--34B & SFT + RLHF & 8--16$\times$ A100/H100 & FSDP + TP=4 for gen \\
70B & RLHF (full) & 32--64$\times$ A100/H100 & Decoupled, FSDP + TP=8 \\
70B & RLHF (LoRA) & 8--16$\times$ A100/H100 & No ref model, LoRA adapters \\
$>$100B & RLHF & 128+$\times$ H100 & 3D parallelism (TP+PP+DP) \\
\bottomrule
\end{tabular}
\end{table}

\begin{intuitionbox}[H100 vs A100: When is the Upgrade Worth It?]
H100 provides:

\begin{itemize}
  \item $\sim$1.6$\times$ peak FLOPS (989 vs 624 TFLOPS for BF16 with sparsity; 495 vs 312 without sparsity)
  \item $\sim$2$\times$ memory bandwidth (3.35 vs 2.0 TB/s)
  \item FP8 support (additional 2$\times$ for inference)
  \item NVLink 4.0 (900 vs 600 GB/s)
\end{itemize}

\textbf{For training}: $\sim$1.8--2.2$\times$ faster end-to-end (FP8 support and higher bandwidth amplify the raw FLOPS advantage).

\textbf{For generation}: $\sim$1.7$\times$ faster (bandwidth-bound, so 2$\times$ BW $\approx$ 1.7$\times$ throughput with overhead).

\textbf{Cost-performance}: At 1.5$\times$ the price, H100 is almost always better value for training. For inference-only (generation clusters), A100 at spot pricing can be more cost-effective.
\end{intuitionbox}

\section{Optimizer Configuration for RL Training}
\label{sec:rl-optimizer-config}

RL training (PPO, GRPO, DPO) imposes unique demands on the optimizer compared to pretraining or SFT. The loss landscape is non-stationary (the policy changes what data is generated), gradients are noisier (reward signal variance), and training is more prone to catastrophic forgetting or reward hacking. This section consolidates RL-specific optimizer guidance, using AdamW~\cite{loshchilov2019adamw} as the default optimizer.

\subsection{Why RL Requires Different Optimizer Settings}
\label{why-rl-requires-different-optimizer-settings}

\begin{keybox}[RL vs. SFT Optimization – Key Differences]
\begin{itemize}
  \item \textbf{Non-stationary data distribution}: unlike SFT where the dataset is fixed, RL generates new rollouts each iteration---the data distribution shifts with the policy.
  \item \textbf{High gradient variance}: reward signals are sparse and noisy; gradients have much higher variance than cross-entropy on curated data.
  \item \textbf{Smaller updates required}: the policy must stay close to the reference model (KL constraint), so learning rates are 10--100$\times$ smaller than SFT.
  \item \textbf{No weight decay}: regularization comes from the KL penalty, not weight decay. Adding WD on top can fight the KL constraint.
  \item \textbf{Shorter warmup}: RL starts from a converged SFT checkpoint---the optimizer state needs minimal warmup.
\end{itemize}
\end{keybox}

\subsection{Recommended Hyperparameters by RL Method}
\label{recommended-hyperparameters-by-rl-method}

\begin{table}[ht!]
\centering
\caption{Optimizer settings for RL training phases. All use $\beta_1=0.9$, $\beta_2=0.95$, $\epsilon=10^{-8}$, \texttt{max\_grad\_norm}=1.0, BF16.}
{\footnotesize
\begin{tabular}{@{}llllll@{}}
\toprule
\textbf{Method} & \textbf{Optimizer} & \textbf{LR} & \textbf{WD} & \textbf{Warmup} & \textbf{Schedule} \\
\midrule
DPO & AdamW & $5\text{e-}7$ & 0.0 & 50 steps & Constant or Linear \\
PPO (policy) & AdamW & $1\text{e-}6$ & 0.0 & 20 steps & Constant \\
PPO (critic) & AdamW & $1\text{e-}6$ & 0.0 & 20 steps & Constant \\
GRPO & AdamW & $1\text{e-}6$ & 0.0 & 20 steps & Constant \\
\bottomrule
\end{tabular}
}
\end{table}

\begin{intuitionbox}[Why Constant Schedule for RL?]
Cosine and linear-decay schedules assume a fixed training horizon and monotonically decreasing loss. RL training has neither: reward may plateau, spike, or oscillate unpredictably. A constant LR (after brief warmup) keeps the optimizer responsive throughout training. If you must decay, use a very gentle linear schedule with a high minimum LR ratio ($\geq 0.5$).
\end{intuitionbox}

\subsection{Beta-2 = 0.95 for RL: Faster Adaptation}
\label{beta-2-0.95-for-rl-faster-adaptation}

The default Adam $\beta_2 = 0.999$ gives a very long memory for the second moment ($\sim$1000-step effective window). In RL training, the loss landscape changes rapidly as the policy evolves---the gradient variance from 1000 steps ago is irrelevant. Using $\beta_2 = 0.95$ shortens the window to $\sim$20 steps, making the adaptive learning rate respond quickly to changing gradient statistics.

\begin{warningbox}[When beta2 = 0.95 Hurts]
For very small batch sizes (e.g., batch=1 in online RL), $\beta_2 = 0.95$ can make the second moment estimates too noisy. In this regime, use $\beta_2 = 0.99$ as a compromise, or increase the effective batch size via gradient accumulation.
\end{warningbox}

\subsection{Mixed Precision for RL: FP32 Master Weights Are Critical}
\label{mixed-precision-for-rl-fp32-master-weights-are-critical}

RL training is particularly sensitive to numerical precision:

\begin{itemize}
  \item Gradients are noisier---small updates must accumulate accurately over many steps
  \item Learning rates are very small ($10^{-6}$--$10^{-7}$), making $\Delta\theta \ll \theta$
  \item BF16 mantissa (7 bits $\approx$ 0.8\% relative precision) cannot represent updates of magnitude $10^{-6}$ relative to weights of magnitude $10^{0}$
\end{itemize}

\textbf{Always use FP32 master weights for RL training.} BF16-only training (no FP32 copy) reliably causes reward collapse in PPO/GRPO after 100--500 steps.

\subsection{Gradient Clipping is Critical for RL}
\label{gradient-clipping-is-critical-for-rl}

In PPO and GRPO, the reward signal can be highly variable, especially early in training. A single bad batch can produce gradients with norm $>100$, which would completely destroy the model weights. \texttt{max\_grad\_norm=1.0} is the standard setting. For SFT, clipping is less critical but still recommended.

\begin{warningbox}[Never Disable Gradient Clipping for RL]
Unlike SFT where gradient norms are typically stable (0.1--1.0 range), RL gradients are spiky because: (1)~reward variance propagates through the policy gradient, (2)~rare high-reward trajectories create outsized updates, and (3)~the KL penalty term can produce large gradients when the policy drifts. A single unclipped step with $\|\nabla\| > 50$ can undo hundreds of training steps.
\end{warningbox}

\subsection{Diagnosing RL Training Instability}
\label{diagnosing-rl-training-instability}

\begin{examplebox}[Red Flags and Fixes for RL Optimization]
{\small
\begin{tabular}{@{}lp{9.5cm}@{}}
\toprule
\textbf{Symptom} & \textbf{Likely Cause and Fix} \\
\midrule
Reward improves then collapses & LR too high or KL coefficient too low. Reduce LR by 2--5$\times$ or increase $\beta_\text{KL}$. \\
Gradient norm constantly at clip threshold & Updates too aggressive. Reduce LR (clipping means you’re losing gradient direction info every step). \\
KL divergence explodes ($>$15 nats) & LR too high. Reduce by 10$\times$ or add adaptive KL penalty. \\
Reward stuck at baseline & LR too low, or reward model has low signal. Try 2--5$\times$ higher LR. Check reward model calibration. \\
Loss NaN after 100+ steps & FP32 master weights missing, or grad norm overflow. Enable FP32 master weights; verify BF16 mode. \\
\bottomrule
\end{tabular}
}
\end{examplebox}

\subsection{HuggingFace TRL Configuration for RL}
\label{huggingface-trl-configuration-for-rl}

The TRL library~\cite{vonwerra2022trl} provides production-ready implementations of PPO, DPO, and other RL methods for LLMs.

\begin{lstlisting}[style=pythonstyle, caption={Complete PPO and DPO optimizer configuration using TRL.}]
from trl import PPOConfig, PPOTrainer, DPOConfig, DPOTrainer

# --- PPO Configuration ---
ppo_config = PPOConfig(
    # Optimizer (AdamW with RL-specific settings)
    learning_rate=1e-6,           # 10-100x smaller than SFT
    
    # PPO-specific
    ppo_epochs=4,                 # mini-batch updates per rollout
    mini_batch_size=16,
    batch_size=64,                # rollout batch size
    
    # Gradient control
    max_grad_norm=1.0,
    
    # KL penalty (replaces weight decay as regularizer)
    init_kl_coef=0.2,            # initial KL penalty coefficient
    adap_kl_ctrl=True,           # adaptive KL targeting
    target_kl=6.0,               # target KL divergence
    
    # Mixed precision
    bf16=True,                   # BF16 compute, FP32 master weights
)

ppo_trainer = PPOTrainer(
    model=model,
    ref_model=ref_model,
    config=ppo_config,
    tokenizer=tokenizer,
    dataset=dataset,
)

# --- DPO Configuration ---
dpo_config = DPOConfig(
    output_dir="./dpo_output",
    
    # Optimizer
    learning_rate=5e-7,           # even smaller than PPO
    optim="adamw_torch",
    adam_beta1=0.9,
    adam_beta2=0.95,              # shorter memory for RL
    weight_decay=0.0,            # no WD -- KL provides regularization
    
    # Schedule
    lr_scheduler_type="constant_with_warmup",
    warmup_steps=50,
    
    # Gradient control
    max_grad_norm=1.0,
    
    # DPO-specific
    beta=0.1,                    # KL constraint strength
    loss_type="sigmoid",         # standard DPO loss
    
    # Mixed precision
    bf16=True,
    
    # Training
    num_train_epochs=1,          # DPO typically 1 epoch
    per_device_train_batch_size=4,
    gradient_accumulation_steps=8,
)

dpo_trainer = DPOTrainer(
    model=model,
    ref_model=ref_model,
    args=dpo_config,
    train_dataset=dataset,
    tokenizer=tokenizer,
)
dpo_trainer.train()
\end{lstlisting}

\subsection{MoE Considerations for RL Training}
\label{moe-considerations-for-rl-training}

\begin{intuitionbox}[MoE for RLHF]
Mixture-of-Experts (MoE) models~\cite{fedus2022switch} are increasingly used in RLHF:

\begin{itemize}
  \item \textbf{Advantage}: 3--4$\times$ more capacity at same compute cost. Better for reward models (more capacity to judge).
  \item \textbf{Challenge}: Expert parallelism requires all-to-all communication (tokens routed across GPUs). This conflicts with pipeline parallelism.
  \item \textbf{GRPO with MoE}: Works well since generation cost is dominated by active params (not total params).
  \item \textbf{LoRA for MoE}: Can apply LoRA to router + shared layers only, or to all experts (expensive).
\end{itemize}
\end{intuitionbox}

\begin{intuitionbox}[The RL Optimizer Mantra]
For RL fine-tuning: \textbf{small LR, no weight decay, constant schedule, FP32 master weights, aggressive clipping}. Let the KL penalty handle regularization---the optimizer’s job is just to follow the policy gradient without overshooting.
\end{intuitionbox}

